% !TeX program = lualatex
% !TeX encoding = utf8
% !TeX spellcheck = uk_UA

\documentclass[]{beamer}
\usetheme{PlasmaPhysics}
\usepackage{PlasmaPhysics}
\usepackage{mathtools}
\setlength{\abovedisplayskip}{3pt}
\setlength{\belowdisplayskip}{3pt}
\usetikzlibrary{decorations.pathreplacing}
\newtcolorbox{mysidebyside}[1][]{
    notitle, sidebyside, blanker,
    sidebyside align=center seam,
    halign=center,
    lefthand width=.5\textwidth,
    #1
}



\graphicspath{{pictures/}}
\title[Лекції з фізики плазми]{\bfseries\huge Адіабатичні інваріанти та магнітні пастки}
\subtitle{\bfseries Лекції з фізики плазми}
\author{Пономаренко С. М.}
\date{}
% https://www.youtube.com/watch?v=zv9Y7YmHyBk
\def\stext#1{\text{
        \makebox[0pt]{
            \tikz[baseline] {\node[text width=50pt, align=center, execute at begin node=\setlength{\baselineskip}{-0.5ex}, blue] {#1};}
        }
}}


\begin{document}




%===============================================================================
\usebackgroundtemplate{
	\tikz\node[opacity=0.1]{\includegraphics[width=\paperwidth,height=\paperheight]{background}};%
}


\begin{frame}
	\thispagestyle{empty}
	\titlepage
\end{frame}



% ============================== Слайд ## ===================================
\begin{frame}{Рух частинки в неоднорідному полі $\nabla B \parallel \vec{B}$}{Магнітне дзеркало (пробкотрон)}
\begin{block}{}

\end{block}

\end{frame}
% ===========================================================================





\end{document}

